상호작용은 다음과 같이 진행된다.

먼저 정수 $n$을 입력받는다.

\begin{center}
$2 \le n \le 1000$
\end{center}

보물은 정확히 두 개이며, 서로 다른 두 위치에 존재한다.

그 후, 최대 $25$번까지 질의할 수 있다.
질의를 하고 싶다면 다음 형식으로 한 줄을 출력해야 한다.
\begin{center}
\texttt{? } $k$ $i_1$ $i_2$ \dots $i_k$
\end{center}

여기서 다음 조건을 만족해야 한다.
\begin{itemize}
\item $1 \le k \le n$
\item $1 \le i_j \le n$
\item $i_1, i_2, \dots, i_k$는 모두 서로 달라야 한다
\end{itemize}

질의를 출력한 뒤에는 반드시 출력 버퍼를 flush해야 한다.
그 후 인터랙터로부터 다음 중 하나의 값을 입력받는다.
\begin{itemize}
\item \texttt{1}: 질의한 위치들 중 적어도 하나에 보물이 있는 경우
\item \texttt{0}: 질의한 모든 위치에 보물이 없는 경우
\end{itemize}

정답을 출력하고 싶다면 다음 형식으로 한 줄을 출력해야 한다.
\begin{center}
\texttt{! } $x$ $y$
\end{center}

여기서 $x, y$는 보물이 있는 서로 다른 두 위치여야 한다.

정답을 출력한 뒤에도 반드시 출력 버퍼를 flush한 후 즉시 프로그램을 종료해야 한다.

질문 형식이 잘못되었거나, 범위를 벗어났거나, 질문 횟수 제한을 초과하면 \textbf{틀렸습니다}를 받는다.

최종 정답의 형식이 잘못되었거나, 출력한 두 위치가 정답과 다르면 \textbf{틀렸습니다}를 받는다.

출력을 즉시 채점 프로그램에 전달하기 위해, 질문이나 정답을 출력한 뒤에는 반드시 출력 버퍼를 비워야 한다.
이를 지키지 않아 상호작용이 정상적으로 진행되지 않으면 \textbf{런타임에러} 나 \textbf{시간초과}와 같은 판정을 받을 수 있다.

언어별로 출력 버퍼를 비우는 방법의 예시는 다음과 같다.

\begin{itemize}
    \item \texttt{C}: \texttt{fflush(stdout);}
    \item \texttt{C++}: \texttt{std::cout.flush();}
    \item \texttt{Java}: \texttt{System.out.flush();}
    \item \texttt{Python}: \texttt{stdout.flush()}
\end{itemize}

그 외의 언어를 사용하는 경우에는 해당 언어의 문서를 참고하라.

정답을 출력한 뒤 프로그램은 즉시 종료해야 한다.